\section{Restitution radiale orientée et transport symplectique}
\subsection{Récupération de Planck avec orientation conservée}

Avec les demi-axes étiquetés de l’ellipse, nous définissons
\(R_\star=\sqrt{b_S/a_S}\). Alors
\(q=R_\star^4=e^{-2\eta_{\rm el}}=(1-\xi)/(1+\xi)\),
où \(\xi=C^*/A\) et \(\eta_{\rm el}=\operatorname{artanh}\xi\).
Ce rapport radial conserve les marques des deux axes.

\begin{theorem}[Récupérateur radial de la section de retour]
\label{rec:thm:accion}
Dans la coordinatisation positive, si \(q=R_\star^4\), alors
\begin{align}
 \eta_{\rm ret}&=\alpha\frac{499-501q}{1+q},\label{rec:eq:eta-radio}\\
 \hbar_{\rm ret}&=\alpha\frac{499-501q}{1+q}\,10^{-34}\mathcal U_S,
 \qquad\hbar_{\rm pre}=H_5\,10^{-34}\mathcal U_S,\label{rec:eq:planck-radio}\\
 R_{\rm act}&=\frac{\hbar_{\rm pre}}{\hbar_{\rm ret}}
 =\frac{H_5(1+q)}{\alpha(499-501q)}.\label{rec:eq:ratio-radio}
\end{align}
La branche \(\eta_{\rm ret}>0\) équivaut à \(q<499/501\), à l’intérieur de \(q>0\).

\end{theorem}
\begin{proof}
Les identités angulaires donnent \(\eta_{\rm ret}=\alpha(500\xi-1)\). On substitue \(\xi=(1-q)/(1+q)\) ; le numérateur est \(499-501q\). Le dénominateur et \(\alpha\) sont positifs. La réalisation dimensionnelle et le quotient sont ceux des sections d'action préalablement construites. La décade \(-34\) et l'unité \(\mathcal U_S\) conservent leur dérivation déterminantale antérieure.
\end{proof}

Cette formule explicite la section de Planck dans le contre-angle. En termes de grandeurs d’action, la relation angulaire est
\[
 C^*=2\left(\frac{\hbar_{\rm ret}}{10^{-34}\mathcal U_S}+\alpha\right).
\]
La division par la base dimensionnelle extrait le coefficient adimensionnel ; elle n’ajoute pas une grandeur d’action à une constante sans unités.


Pour l’orientation angulaire \(\varepsilon\in\{-1,+1\}\), nous définissons
\[
 C^*_{\varepsilon}=2\varepsilon(\eta_{\rm ret}+\alpha),\quad
 \xi_\varepsilon=C^*_{\varepsilon}/A,\quad
 q_\varepsilon=\frac{1-\xi_\varepsilon}{1+\xi_\varepsilon}>0.
\]
L’opérateur de récupération sur les deux coordinatisations est
\begin{equation}
 \eta_{\rm ret}=\alpha\left(500\varepsilon
        \frac{1-q_\varepsilon}{1+q_\varepsilon}-1\right).
 \label{rec:eq:eta-orientada}
\end{equation}

\begin{corollary}[Invariance du récupérateur conjugué
]
La transformation \((\varepsilon,q)\mapsto(-\varepsilon,q^{-1})\) conserve \(\eta_{\rm ret}\), \(\hbar_{\rm ret}\) et \(R_{\rm act}\).

\end{corollary}
\begin{proof}
Le rapport \((1-q)/(1+q)\) change de signe lorsqu’on inverse \(q\). Le produit par \(\varepsilon\) reste invariant. On applique \eqref{rec:eq:eta-orientada} puis les représentations d’action.
\end{proof}

Si l’orientation est écartée et que \(\varepsilon=+1\) est indûment maintenu, la même expression produit \(\eta_{\rm ret}(q^{-1})=-\eta_{\rm ret}(q)-2\alpha\). Le registre distingue donc deux opérations algébriquement différentes. L’orientation angulaire \(\varepsilon\) ne s’identifie pas non plus sans autre précision au signe symplectique \(\sigma\) : l’échange des axes est antisymplectique, tandis que la rotation qui les échange préserve la forme d’aire.


\subsection{Entrelacement de phase et d'anisotropie}

Soient
\[
 S_\eta=\operatorname{diag}(e^{\eta/2},e^{-\eta/2}),\quad
 J_0=\begin{pmatrix}0&-1\\1&0\end{pmatrix},\quad
 J_\eta=S_\eta J_0S_\eta^{-1}
       =\begin{pmatrix}0&-e^\eta\\e^{-\eta}&0\end{pmatrix}.
\]
On a \(\det S_\eta=1\), \(J_\eta^2=-I\) et \(C_\eta(\theta)=e^{\theta J_\eta}=S_\eta e^{\theta J_0}S_\eta^{-1}\). La loi \(S_{\eta+\delta}=S_\delta S_\eta\) prouve l'entrelacement exact
\begin{equation}
 C_{\eta+\delta}(\theta)S_\delta=S_\delta C_\eta(\theta).
 \label{rec:eq:entrelazamiento}
\end{equation}
La transformation préserve \(dQ\wedge dP\) et transporte l’évolution de phase vers la nouvelle ellipse. Avec la normalisation de \eqref{eq:ii-elipse-area-fase}, l’action balayée est \(\hbar\theta\), ou \(\sigma\hbar\theta\) dans le feuillet orienté. La phase \(\theta\), l’anisotropie \(\eta_{\rm el}\) et le paramètre \(t\) du flux de \(K\) conservent ainsi des fonctions différentes, reliées par des opérateurs explicites.
